% GENERATED FILE. Edit canonical lessons, then run npm run generate:books.
% canonical-source-hash: fnv1a64:051690b91ad61504
% canonical-lessons: HI-C438-jara, HI-C438-ola, HI-C438-sit, HI-C438-himpat, HI-C438-orhna

\chapter{Winter Cold}
\label{ch:hi-winter-cold}
\hlchaptermodality{438}

\begin{chapteropening}
\textbf{By the end of this chapter:} \emph{I can talk about winter, hail and a cold wave, snowfall in the hills and wrapping up warm.}

Complete the last lesson of chapter 438: I can talk about winter, hail and a cold wave, snowfall in the hills and wrapping up warm.
\end{chapteropening}

\section[\texorpdfstring{\hi{जाड़ा} (jāṛā)}{जाड़ा (jāṛā)}]{\hi{जाड़ा} (jāṛā) — winter; the cold}

\label{lesson:HI-C438-jara}



\begin{quote}
Before the new one: say the Hindi for to spread over, then the Hindi for to sweat.
\end{quote}

\subsection*{You'll want to know: \hi{जाड़ा}}
\textbf{\hi{जाड़ा}} — \emph{jāṛā} — \textquotedblleft{}winter; the cold\textquotedblright{}.

\begin{quote}
\hi{जाड़े} \hi{में} \hi{यहाँ} \hi{बहुत} \hi{ठंड} \hi{पड़ती} \hi{है।}
\end{quote}

\begin{quote}
\emph{jāṛe meṁ yahā̃ bahut ṭhaṇḍ paṛtī hai.}
\end{quote}

\begin{quote}
It gets very cold here in winter.
\end{quote}

Masculine.

\begin{cousinweb}
\textbf{\hi{जाड़ा}} (\emph{jāṛā}) is a tadbhava of Sanskrit \emph{jaḍa}, \textquotedblleft{}dull, numb\textquotedblright{} and also \textquotedblleft{}cold\textquotedblright{} (CDIAL 5090), which lists Hindi \emph{jaṛā}, \textquotedblleft{}cold, frost, winter\textquotedblright{}.
\end{cousinweb}

\begin{grammarlens}[title={What you've built}]
One more word for getting around.
\end{grammarlens}

\subsection*{Your turn}
\begin{itemize}
\raggedright
  \item \emph{Say it:} \emph{jāṛā}
  \item \emph{Say it:} \emph{jāṛā}, once more
  \item \emph{Say it:} \emph{jāṛe meṁ yahā̃ bahut ṭhaṇḍ paṛtī hai.} — It gets very cold here in winter.
  \item \emph{Recall:} say the Hindi for to spread over, then the Hindi for to sweat, then say \emph{jāṛā} again
\end{itemize}

\begin{tcolorbox}[breakable,colback=teal!4,colframe=teal!35!black,title={Before you move on}]
What does \hi{जाड़ा} mean? (\textquotedblleft{}Winter; the cold\textquotedblright{}.) Say it once more.
\end{tcolorbox}
\section[\texorpdfstring{\hi{ओला} (olā)}{ओला (olā)}]{\hi{ओला} (olā) — a hailstone}

\label{lesson:HI-C438-ola}



\begin{quote}
Before the new one: say the Hindi for to sweat, then the Hindi for winter.
\end{quote}

\subsection*{You'll want to know: \hi{ओला}}
\textbf{\hi{ओला}} — \emph{olā} — \textquotedblleft{}a hailstone\textquotedblright{}.

\begin{quote}
\hi{कल} \hi{रात} \hi{ओले} \hi{गिरे।}
\end{quote}

\begin{quote}
\emph{kal rāt ole gire.}
\end{quote}

\begin{quote}
It hailed last night.
\end{quote}

Masculine.

Usually plural: \hi{ओले} \hi{गिरे}, \hi{ओले} \hi{पड़े}, it hailed.

\begin{cousinweb}
\textbf{\hi{ओला}} (\emph{olā}) is a tadbhava of Sanskrit \emph{upala}, \textquotedblleft{}a stone, a rock\textquotedblright{} (CDIAL 2223), which lists Hindi \emph{olā}, \textquotedblleft{}a hailstone\textquotedblright{}: a stone that falls from the sky.
\end{cousinweb}

\begin{grammarlens}[title={What you've built}]
One more word for getting around.
\end{grammarlens}

\subsection*{Your turn}
\begin{itemize}
\raggedright
  \item \emph{Say it:} \emph{olā}
  \item \emph{Say it:} \emph{olā}, once more
  \item \emph{Say it:} \emph{kal rāt ole gire.} — It hailed last night.
  \item \emph{Recall:} say the Hindi for to sweat, then the Hindi for winter, then say \emph{olā} again
\end{itemize}

\begin{tcolorbox}[breakable,colback=teal!4,colframe=teal!35!black,title={Before you move on}]
What does \hi{ओला} mean? (\textquotedblleft{}A hailstone\textquotedblright{}.) Say it once more.
\end{tcolorbox}
\section[\texorpdfstring{\hi{शीत} (śīt)}{शीत (śīt)}]{\hi{शीत} (śīt) — cold; the cold (formal)}

\label{lesson:HI-C438-sit}



\begin{quote}
Before the new one: say the Hindi for winter, then the Hindi for a hailstone.
\end{quote}

\subsection*{You'll want to know: \hi{शीत}}
\textbf{\hi{शीत}} — \emph{śīt} — \textquotedblleft{}cold; the cold (formal)\textquotedblright{}.

\begin{quote}
\hi{शीत} \hi{लहर} \hi{की} \hi{वजह} \hi{से} \hi{स्कूल} \hi{बंद} \hi{हैं।}
\end{quote}

\begin{quote}
\emph{śīt lahar kī vajah se skūl band haiṁ.}
\end{quote}

\begin{quote}
The schools are closed because of the cold wave.
\end{quote}

Masculine.

\begin{cousinweb}
\textbf{\hi{शीत}} (\emph{śīt}) is a tatsama, Sanskrit \emph{śīta}, \textquotedblleft{}cold, chilly\textquotedblright{}, and as a noun \textquotedblleft{}cold weather\textquotedblright{} (Monier-Williams); the weather news speaks of \hi{शीत} \hi{लहर}, a cold wave.
\end{cousinweb}

\begin{grammarlens}[title={What you've built}]
One more word for getting around.
\end{grammarlens}

\subsection*{Your turn}
\begin{itemize}
\raggedright
  \item \emph{Say it:} \emph{śīt}
  \item \emph{Say it:} \emph{śīt}, once more
  \item \emph{Say it:} \emph{śīt lahar kī vajah se skūl band haiṁ.} — The schools are closed because of the cold wave.
  \item \emph{Recall:} say the Hindi for winter, then the Hindi for a hailstone, then say \emph{śīt} again
\end{itemize}

\begin{tcolorbox}[breakable,colback=teal!4,colframe=teal!35!black,title={Before you move on}]
What does \hi{शीत} mean? (\textquotedblleft{}Cold; the cold (formal)\textquotedblright{}.) Say it once more.
\end{tcolorbox}
\section[\texorpdfstring{\hi{हिमपात} (himpāt)}{हिमपात (himpāt)}]{\hi{हिमपात} (himpāt) — snowfall}

\label{lesson:HI-C438-himpat}



\begin{quote}
Before the new one: say the Hindi for a hailstone, then the Hindi for cold (formal).
\end{quote}

\subsection*{You'll want to know: \hi{हिमपात}}
\textbf{\hi{हिमपात}} — \emph{himpāt} — \textquotedblleft{}snowfall\textquotedblright{}.

\begin{quote}
\hi{पहाड़ों} \hi{में} \hi{कल} \hi{रात} \hi{भारी} \hi{हिमपात} \hi{हुआ।}
\end{quote}

\begin{quote}
\emph{pahāṛoṁ meṁ kal rāt bhārī himpāt huā.}
\end{quote}

\begin{quote}
There was heavy snowfall in the hills last night.
\end{quote}

Masculine.

\begin{cousinweb}
\textbf{\hi{हिमपात}} (\emph{himpāt}) is a tatsama, Sanskrit \emph{himapāta}, \textquotedblleft{}a fall of snow\textquotedblright{} (Monier-Williams), from \emph{hima}, \textquotedblleft{}snow\textquotedblright{}, as in \hi{हिमालय}, and \emph{pāta}, \textquotedblleft{}a fall\textquotedblright{}.
\end{cousinweb}

\begin{grammarlens}[title={What you've built}]
One more word for getting around.
\end{grammarlens}

\subsection*{Your turn}
\begin{itemize}
\raggedright
  \item \emph{Say it:} \emph{himpāt}
  \item \emph{Say it:} \emph{himpāt}, once more
  \item \emph{Say it:} \emph{pahāṛoṁ meṁ kal rāt bhārī himpāt huā.} — There was heavy snowfall in the hills last night.
  \item \emph{Recall:} say the Hindi for a hailstone, then the Hindi for cold (formal), then say \emph{himpāt} again
\end{itemize}

\begin{tcolorbox}[breakable,colback=teal!4,colframe=teal!35!black,title={Before you move on}]
What does \hi{हिमपात} mean? (\textquotedblleft{}Snowfall\textquotedblright{}.) Say it once more.
\end{tcolorbox}
\section[\texorpdfstring{\hi{ओढ़ना} (oṛhnā)}{ओढ़ना (oṛhnā)}]{\hi{ओढ़ना} (oṛhnā) — to wrap oneself in (a shawl or a blanket)}

\label{lesson:HI-C438-orhna}



\begin{quote}
Before the new one: say the Hindi for cold (formal), then the Hindi for snowfall.
\end{quote}

\subsection*{You'll want to know: \hi{ओढ़ना}}
\textbf{\hi{ओढ़ना}} — \emph{oṛhnā} — \textquotedblleft{}to wrap oneself in (a shawl or a blanket)\textquotedblright{}.

\begin{quote}
\hi{दादी} \hi{ने} \hi{कंबल} \hi{ओढ़ा।}
\end{quote}

\begin{quote}
\emph{dādī ne kambal oṛhā.}
\end{quote}

\begin{quote}
Grandmother wrapped herself in a blanket.
\end{quote}

Past, with \hi{ने} (he or she …): \hi{उसने} \hi{चादर} \hi{ओढ़ी}, \emph{usne cādar oṛhī}.

\begin{cousinweb}
\textbf{\hi{ओढ़ना}} (\emph{oṛhnā}) goes back to \emph{*ōḍḍh}, \textquotedblleft{}put on, wear, cover\textquotedblright{}, a form CDIAL reconstructs and thinks may come from Dravidian (CDIAL 2547); it lists Hindi \emph{oṛhnā}, \textquotedblleft{}to cover\textquotedblright{}.
\end{cousinweb}

\begin{grammarlens}[title={What you've built}]
One more word for getting around.
\end{grammarlens}

\subsection*{Your turn}
\begin{itemize}
\raggedright
  \item \emph{Say it:} \emph{oṛhnā}
  \item \emph{Say it:} \emph{oṛhnā}, once more
  \item \emph{Say it:} \emph{dādī ne kambal oṛhā.} — Grandmother wrapped herself in a blanket.
  \item \emph{Recall:} say the Hindi for cold (formal), then the Hindi for snowfall, then say \emph{oṛhnā} again
\end{itemize}

\begin{tcolorbox}[breakable,colback=teal!4,colframe=teal!35!black,title={Before you move on}]
What does \hi{ओढ़ना} mean? (\textquotedblleft{}To wrap oneself in (a shawl or a blanket)\textquotedblright{}.) Say it once more.
\end{tcolorbox}
